% Authored native TikZ addition; embedded by native_figures.py.

% Top panels: every arrow denotes an electron transition, not hole motion.
\foreach \xx/\letter/\title in {0/a/电子捕获,3.5/b/电子发射,7/c/空穴捕获,10.5/d/空穴发射}{
 \begin{scope}[xshift=\xx cm]
  \node[font=\small\bfseries] at(1.05,2.9) {(\letter) {\cjkfont \title}};
  \draw[line width=.65pt] (0,2.2)--(2.15,2.2) node[right,font=\small] {$E_C$};
  \draw[dashed,line width=.65pt] (.2,1.1)--(1.95,1.1) node[right,font=\small] {$E_t$};
  \draw[line width=.65pt] (0,0)--(2.15,0) node[right,font=\small] {$E_V$};
 \end{scope}
}
\draw[-{Stealth},line width=.6pt](-.65,0)--(-.65,2.25);
\node[rotate=90,font=\small]at(-1,1.15){{\cjkfont 电子能量}};
\draw[-{Stealth},tolblue,line width=1.15pt](1.05,2.15)--(1.05,1.16);
\draw[-{Stealth},tolblue,line width=1.15pt](4.55,1.16)--(4.55,2.15);
\draw[-{Stealth},tolorange,line width=1.15pt](8.05,1.04)--(8.05,.06);
\draw[-{Stealth},tolorange,line width=1.15pt](11.55,.06)--(11.55,1.04);
\node[font=\small] at(1.05,-.50) {$0\to1$};
\node[font=\small] at(4.55,-.50) {$1\to0$};
\node[font=\small] at(8.05,-.50) {$1\to0$};
\node[font=\small] at(11.55,-.50) {$0\to1$};
\node[font=\small] at(1.05,-1.05) {$r_{nc}=N_t c_n n(1-f)$};
\node[font=\small] at(4.55,-1.05) {$r_{ne}=N_t e_n f$};
\node[font=\small] at(8.05,-1.05) {$r_{pc}=N_t c_p p f$};
\node[font=\small] at(11.55,-1.05) {$r_{pe}=N_t e_p(1-f)$};
\node[font=\small,align=center] at(1.05,-1.63) {{\cjkfont 移走一个}\ {\cjkfont 传导电子}};
\node[font=\small,align=center] at(4.55,-1.63) {{\cjkfont 增加一个}\ {\cjkfont 传导电子}};
\node[font=\small,align=center] at(8.05,-1.63) {{\cjkfont 填补已有价态空穴}};
\node[font=\small,align=center] at(11.55,-1.63) {{\cjkfont 留下一个价态空穴}};
% Bottom arrows denote the center's occupation changes, not energy or space.
\node[box,minimum width=2.8cm,minimum height=1.15cm] (empty) at(1.7,-3.55) {$\circ\quad 0$ {\cjkfont 空态}\\{\cjkfont 概率} $1-f$};
\node[box,draw=tolgreen,minimum width=2.8cm,minimum height=1.15cm] (filled) at(10.9,-3.55) {$\bullet\quad 1$ {\cjkfont 电子占据}\\{\cjkfont 概率} $f$};
\draw[-{Stealth},line width=.85pt,draw=tolblue] (3.2,-3.24)--node[above,font=\small] {{\cjkfont 电子捕获} $c_n n$ {\cjkfont ＋空穴发射} $e_p$}(9.4,-3.24);
\draw[-{Stealth},line width=.85pt,draw=tolorange] (9.4,-3.85)--node[below,font=\small] {{\cjkfont 电子发射} $e_n$ {\cjkfont ＋空穴捕获} $c_p p$}(3.2,-3.85);
